\section{Density Zeros and Support Rigidity}

A smooth density can vanish outside an interval. Contract just one
coordinate, and part of that interval becomes inaccessible to the image
law. The change has rank one and can be arbitrarily small, yet equivalence
is lost. This obstruction is already present in finite dimension; no
accumulation of errors is needed.

We separate this support obstruction from the Hellinger cost. The local
geometry still measures small matrix perturbations, and absolute continuity
still forces Hilbert--Schmidt proximity to a coordinate symmetry. Within
that neighbourhood, the singular part is exactly the mass sent into zeros
of the density. Requiring the same avoidance for the inverse gives
equivalence. For bounded and one-sided laws, these support constraints
determine the operator outright.

\subsection{Separating the mixing cost from the support obstruction}

The path argument differentiates the square-root density on the whole
line. It continues to work when that function reaches zero continuously
enough to remain Sobolev. We assume
\begin{equation}
 \begin{gathered}
 f\geq0,\quad \int f=1,\quad
 g=\sqrt f\in W^{1,2}(\R),\quad xg'\in L^2(\R),\\
 \E X=0,\quad \E X^2=1,\quad f\text{ non-Gaussian}.
 \end{gathered}
 \label{eq:zero-set-row-assumptions}
\end{equation}
The boundary is part of this regularity condition. For example, the uniform
density has zero derivative inside its support, but its square root jumps
at the endpoints. Such a jump contributes a first-order Hellinger loss
under a small rotation, as we will see below. A Sobolev square root has no
boundary jump, and its derivative vanishes almost everywhere on its zero
set. The identity \(sg=-2g'\) therefore survives on the whole line.

\begin{proposition}[Local Fourier coercivity with zeros of the density]
\label{prop:local-coercivity-with-density-zeros}
Under \eqref{eq:zero-set-row-assumptions}, the matrix Fisher operator decomposition,
finite-dimensional Hellinger path bound, and dimension-free local
Frobenius coercivity in
\cref{sec:dimension-free-hellinger-geometry} remain valid.
Write \(\mu=\rho^{\otimes\N},\quad\rho(dx)=f(x)\,dx\).
There are constants \(c,C,\delta>0\), with \(\delta<1\), such that
\begin{equation}
 c\|K\|_{\mathcal S_2}\leq h(\mu,\nu_{I+K})
       \leq C\|K\|_{\mathcal S_2}
 \qquad(\|K\|_{\mathcal S_2}\leq\delta).
 \label{eq:local-hs-bound-with-density-zeros}
\end{equation}
Moreover, for every \(T\in\GL(H)\), the following necessary
condition survives:
\begin{equation}
 \nu_T\ll\mu
 \quad\Longrightarrow\quad
 T-Q\in\mathcal S_2
 \quad\text{for some }Q\in\mathcal P.
 \label{eq:hs-necessity-with-density-zeros}
\end{equation}
For orthogonal transformations one may omit \(xg'\in L^2\),
retaining non-Gaussianity, \(g\in W^{1,2}(\R)\), and normalization to a
centered law with variance one.
\end{proposition}
\begin{proof}
Define \(s=-2g'/g\) on \(\{g>0\}\) and zero elsewhere.  Since Sobolev
derivatives vanish almost everywhere on zero sets, \(sg=-2g'\) globally, and
all cutoff identities---including \(\E s(X)=0\), \(\E[Xs(X)]=1\), and the
formulas for \(J,\kappa\)---remain valid.  The cases \(J=1\) and \(\kappa=0\)
would respectively give \(g'=-xg/2\) and \(2xg'+g=0\); the latter has no
nonzero square-integrable solution.  Thus the finite-dimensional proofs apply.

For \(K\in\mathcal S_2\) with \(I+K\) invertible, put
\(T_n=I+P_nKP_n\).
The Hellinger path bound and invariance give, for large \(m,n\),
\[
 h(\nu_{T_n},\nu_{T_m})
 \leq C\|T_n^{-1}\|_{\mathrm{op}}
                 \|T_n-T_m\|_{\mathcal S_2}.
\]
The square-root densities are therefore Cauchy in \(L^2\) relative to a common
dominating measure.  Their norm-one limit is identified with \(\nu_{I+K}\) by
finite-row convergence, proving \eqref{eq:local-hs-bound-with-density-zeros}.
The dominating measure here includes the parts outside the original
positivity set. Polar approximants and the skew-generator bound
from \cref{prop:orthogonal-sufficiency} give the orthogonal case.

In \cref{lem:hilbert-schmidt-necessity}, the tail-symmetry and covariance
lemmas already require only \(\nu_T\ll\mu\); the other inputs are moment-only
rigidity, the measurable inverse action, and local coercivity for finite-rank
tail perturbations.  The estimate above supplies the last input, so the same
tail-matching proof yields \eqref{eq:hs-necessity-with-density-zeros}.
\end{proof}

To identify absolute continuity, we now ask where the image mass goes.
Finite-coordinate approximants have densities wherever the original product
is positive, and Hellinger convergence preserves this absolute continuity
on that set. All singular mass must therefore lie outside it. For
equivalence, apply the same argument to the inverse.

\begin{theorem}[The linear equivalence group with density zeros]
\label{thm:linear-equivalence-group}
Assume \eqref{eq:zero-set-row-assumptions} and put
$\mathcal U=\{x\in\R:f(x)>0\}$. For every $T\in\GL(H)$,
\begin{equation}
 \nu_T\ll\mu
 \quad\Longleftrightarrow\quad
 \begin{cases}
 T-Q\in\mathcal S_2(H)\text{ for some }Q\in\mathcal P,\\
 \displaystyle
 \mu\!\left\{x:\sum_jT_{ij}x_j\notin\mathcal U\right\}=0
       \text{ for every }i.
 \end{cases}
 \label{eq:absolute-continuity-operator-criterion}
\end{equation}
Consequently,
\[
 \{T\in\GL(H):\nu_T\sim\mu\}
 =\left\{T\in\GL(H):
 \begin{array}{l}
 T-Q\in\mathcal S_2(H)\text{ for some }Q\in\mathcal P,\\
 \displaystyle\mu\!\left\{x:\sum_jT_{ij}x_j\notin\mathcal U\right\}=0,\\
 \displaystyle\mu\!\left\{x:\sum_j(T^{-1})_{ij}x_j\notin\mathcal U\right\}=0
       \quad\text{for every }i
 \end{array}\right\}.
\]
It is a subgroup of $\GL(H)$.  For orthogonal members the criterion holds with
$T^{-1}=T^*$ and without $xg'\in L^2$.

Whenever $T-Q\in\mathcal S_2(H)$ for some
$Q\in\mathcal P$, the Lebesgue decomposition of $\nu_T$ relative to
$\mu$ is its restriction to $\mathcal U^{\N}$ plus its restriction
to the complement: the former is absolutely continuous and the latter
is singular.
\end{theorem}
\begin{proof}
Suppose $T-Q\in\mathcal S_2$ and put $K=TQ^{-1}-I$.  Since $\nu_Q=\mu$,
$\nu_T=\nu_{I+K}$.  For $T_n=I+P_nKP_n$,
\cref{prop:local-coercivity-with-density-zeros} gives
$h(\nu_{T_n},\nu_T)\to0$, hence total-variation convergence.  Each $\nu_{T_n}$ has a
Lebesgue density in its first $n$ coordinates and agrees with $\mu$ on the
tail, so positivity of $\mu_n$ on $\mathcal U^n$ and Fubini give
$\nu_{T_n}|_{\mathcal U^{\N}}\ll\mu$.  Passing to the limit proves the same
for $\nu_T$, while
$\nu_T|_{(\mathcal U^{\N})^c}\perp\mu$.  Consequently
\[
 \nu_T\ll\mu
 \quad\Longleftrightarrow\quad
 \nu_T(\mathcal U^{\N})=1
 \quad\Longleftrightarrow\quad
 \mu\!\left\{x:\sum_jT_{ij}x_j\notin\mathcal U\right\}=0
       \text{ for every }i.
\]
Together with \eqref{eq:hs-necessity-with-density-zeros}, this proves
\eqref{eq:absolute-continuity-operator-criterion}.

If $\nu_T\sim\mu$, the measurable inverse action gives
$\nu_{T^{-1}}\sim\mu$, so the forward and inverse zero-set row conditions are
necessary.  Conversely,
$T^{-1}-Q^{-1}=-T^{-1}(T-Q)Q^{-1}\in\mathcal S_2$, and these conditions with
\eqref{eq:absolute-continuity-operator-criterion} give
$\nu_T,\nu_{T^{-1}}\ll\mu$.  The inverse action then yields
\[
 \mu=(\Phi_T)_\#\nu_{T^{-1}}\ll(\Phi_T)_\#\mu=\nu_T,
\]
and hence equivalence.  Measurable composition gives
$\nu_{ST}=(\Phi_S)_\#\nu_T\sim\nu_S\sim\mu$ and closure under inverses follows
above, proving the group assertion.  In the orthogonal case, use the
approximants from \cref{prop:orthogonal-sufficiency}, which require only
location Fisher information.
\end{proof}

The forward condition prevents the image from entering a forbidden set.
The inverse condition prevents it from leaving behind a region of positive
original mass. These are different requirements, as the rank-one
contraction below illustrates. Since nonzero row sums have densities,
changing \(f\) on a Lebesgue-null set changes neither condition.

\subsection{A rotating square: overlap without equivalence}

The difference between boundary loss and accumulated loss is visible in a
square. A small rotation removes thin triangles from its support. With
uniform density, their mass is proportional to the angle; with a density
vanishing smoothly at the boundary, their effect can be smaller. Start
with \(\rho\) uniform on \([-\sqrt3,\sqrt3]\), let
\(\mu_{\mathrm u,2}=\rho^{\otimes2}\), and consider
\(T=\bigoplus_kO_{\theta_k}\).
Write
\(\delta_k=\min_{m\in\mathbb Z}|\theta_k-m\pi/2|\in[0,\pi/4]\).
For one block the affinity is the relative area of intersection
of a square and its rotation.
Removing the four corner triangles gives, for \(0<\delta\leq\pi/4\),
\[
 \begin{aligned}
 \Aff(\mu_{\mathrm u,2},(O_\delta)_\#\mu_{\mathrm u,2})
 &=1-\frac{(\cos\delta+\sin\delta-1)^2}
                {2\sin\delta\cos\delta},\\
 1-\Aff(\mu_{\mathrm u,2},(O_\delta)_\#\mu_{\mathrm u,2})
 &=\frac{\delta}{2}+O(\delta^2).
 \end{aligned}
\]
At \(\delta=0\), the affinity is one. Each nonzero \(\delta\) leaves
positive area on both sides of the support difference, precluding absolute
continuity in either direction already in that block. Independence makes
the affinity of finitely many pairs the product of their affinities.
Passing to the limit by \eqref{eq:general-marginal-affinity-limit} gives
\[
 \Aff(\mu,\nu_T)
 =\prod_k\Aff(\mu_{\mathrm u,2},(O_{\delta_k})_\#\mu_{\mathrm u,2}).
\]
The expansion above makes this product positive exactly when
\(\sum_k\delta_k<\infty\). Thus
\[
 \begin{array}{c|l}
  \delta_k=0\text{ for every }k&\nu_T=\mu\\
  0<\sum_k\delta_k<\infty&
      \nu_T\not\perp\mu,\quad \nu_T\not\ll\mu,\quad\mu\not\ll\nu_T\\
 \sum_k\delta_k=\infty&\nu_T\perp\mu.
 \end{array}
\]
In particular, \(\theta_k=1/k\) gives \(T-I\in\mathcal S_2\) but
\(\nu_T\perp\mu\).
With \(\theta_k=1/k^2\), the laws instead have positive affinity
and neither is absolutely continuous with respect to the other.

To see the effect of a smoother boundary, let \(f\) be the variance-one rescaling of
\((1-x^2)^r\mathbf1_{\{|x|<1\}}\), \(r>1\), put
\(\mu_{f,2}=(f(x)\,dx)^{\otimes2}\), and set
\(J=4\|(\sqrt f)'\|_2^2>1\).
The skew-generator formula gives
\begin{equation}
 1-\Aff(\mu_{f,2},(O_\theta)_\#\mu_{f,2})
       =\frac{J-1}{4}\theta^2+o(\theta^2).
 \label{eq:positive-density-rotation-affinity}
\end{equation}
Only signed-permutation rotations preserve the square support, but every
rotation retains positive overlap near the origin.  Continuity and
\eqref{eq:positive-density-rotation-affinity} therefore replace
\(\sum_k\delta_k\) by \(\sum_k\delta_k^2\) in the trichotomy, while equivalence
still requires every \(\delta_k=0\).

\subsection{When the support determines the operator}

A bounded coordinate cannot be mixed freely with many independent copies
without enlarging its range. Conditioning on finitely many inputs makes
this observation precise: the width of the resulting sum is the original
width multiplied by the \(\ell^1\)-norm of the coefficients. If both the
map and its inverse preserve the range, both matrices have row norms at
most one. Their inverse identities then force a signed permutation.
For a half-line, a negative coefficient is ruled out by the unbounded tail,
and the same argument leaves only ordinary permutations.

Write \(\rho=\Law(X)\), \(\rho^-=\Law(-X)\), and \(\mu=\rho^{\otimes\N}\),
with \(\E X=0\) and \(\E X^2=1\).

\begin{theorem}[Linear equivalence for bounded and half-line coordinate laws]
\label{thm:bounded-and-half-line-classification}
Let \(T\in\GL(H)\).
\begin{enumerate}
 \item Suppose \(a=\operatorname{ess\,inf}X<0<
 b=\operatorname{ess\,sup}X\) are finite.
 Then \(\nu_T\sim\mu\) forces \(T\) to be a signed permutation.
 The complete equivalence group is as follows:
 \[
 \begin{array}{c|l}
 \text{relation between }\rho^-\text{ and }\rho&\text{allowed coordinate signs}\\ \hline
 \rho^-=\rho&\text{all signs}\\
 \rho^-\sim\rho,\ \rho^-\ne\rho&
                  \text{only finitely many negative signs}\\
 \rho^-\not\sim\rho&\text{no negative signs}.
 \end{array}
 \]
 In particular, \(a+b\ne0\) permits only ordinary permutations.
 The support need not fill an interval, and a density is not required.
 \item Suppose \(\operatorname{ess\,inf}X=a>-\infty\) and
 \(\operatorname{ess\,sup}X=+\infty\).
 Then
 \[
  \nu_T\sim\mu
   \quad\Longleftrightarrow\quad
  T\text{ is an ordinary coordinate permutation}.
 \]
 The same conclusion holds for a law bounded above and unbounded below.
\end{enumerate}
In the second case every equivalence is exact invariance.
\end{theorem}
\begin{proof}
Assume \(\nu_T\sim\mu\), and set \(S=T^{-1}\).
The measurable inverse action also gives \(\nu_S\sim\mu\).
For every row \(v\) of \(T\) or \(S\), the series
\(Z=\sum_jv_jX_j\) converges in \(L^2\), and
\begin{equation}
 \E[Z\mid X_1,\ldots,X_m]=\sum_{j\leq m}v_jX_j.
 \label{eq:finite-row-conditional-expectation}
\end{equation}

In the bounded case \(Z\in[a,b]\) almost surely, so its conditional
expectation does too.  The essential width of the finite sum in
\eqref{eq:finite-row-conditional-expectation} is
\((b-a)\sum_{j\leq m}|v_j|\); hence every row of \(T\) and \(S\) has
\(\ell^1\)-norm at most one.

Let \(A_{ij}=|T_{ij}|\) and \(B_{ij}=|S_{ij}|\).
The \(\ell^2\) row/column Cauchy--Schwarz inequality makes the entries of
\(TS\) and \(ST\) absolutely convergent, so \(AB\ge I\) and \(BA\ge I\)
entrywise.  Tonelli and the row bounds give
\[
 \sum_j(AB)_{ij}
 =\sum_k A_{ik}\sum_jB_{kj}\leq1,
 \qquad
 \sum_j(BA)_{ij}\leq1.
\]
Thus \(AB=BA=I\).  For each \(i\), choose \(k\) with
\(A_{ik}B_{ki}>0\).  The off-diagonal entries of \(AB\) force row \(k\) of
\(B\) to vanish outside column \(i\), and those of \(BA\) force row \(i\)
of \(A\) to vanish outside column \(k\).  This gives a bijective monomial
pattern.  Its weights and reciprocals are at most one by the row bounds, so
every weight is one; hence \(T\) is a signed permutation.

For a negative row sign, scalar equivalence requires
\(\rho^-\sim\rho\), and therefore \(a=-b\).  Each negative sign contributes
\(\Aff(\rho,\rho^-)\) to the product affinity.  This number equals one
exactly when \(\rho^-=\rho\); otherwise Kakutani's theorem
\cite{kakutani1948equivalence} permits only finitely many negative signs.
This proves the three alternatives.

In the lower-bounded, upper-unbounded case, centering and nondegeneracy give
\(a<0\).  Every finite sum in
\eqref{eq:finite-row-conditional-expectation} is bounded below by \(a\), so
no coefficient can be negative.  Its essential infimum is
\(a\sum_{j\le m}v_j\), whence
\[
 a\sum_{j\leq m}v_j\geq a,
 \qquad\text{hence}\qquad \sum_jv_j\leq1.
\]
Thus \(T\) and \(S\) are nonnegative with row sums at most one, and the
preceding inverse-matrix argument makes \(T\) an ordinary permutation.  Such
permutations preserve \(\mu\); reflection handles the upper-bounded case.
\end{proof}

If \(Y\) has a Gamma law with shape \(r_\Gamma>0\) and arbitrary positive scale,
its standardization has endpoints \(-\sqrt{r_\Gamma}\) and \(+\infty\); hence only coordinate
permutations are allowed for every positive shape and scale.  For compact
support, signs are governed by \(\rho\) versus \(\rho^-\): if
\(\rho^-\sim\rho\) but \(\rho^-\ne\rho\), only finitely many negative signs are
allowed, whereas unequal endpoint magnitudes, as for a centered asymmetric
Beta law, preclude them.

These support conclusions also cover densities with abrupt boundaries.
For Gamma laws the zero-extended square root is Sobolev exactly when the
shape exceeds two; for Beta laws both shapes must exceed two. The support
argument treats all positive shapes in the same way.

\paragraph{What a one-coordinate contraction leaves behind.}

Transformations outside the equivalence group can retain overlap: the
rotated squares above give neither direction of absolute continuity, while
a contraction gives exactly one. Return to the contraction that opened
this section. It keeps every output in the original support, but leaves
a region of positive mass unreachable. This is what the inverse row
condition detects.

Let \(f\) be the rescaling to variance one of
\(C(1-x^2)^2\mathbf 1_{\{|x|<1\}}\), and write its support as
\([-a,a]\).
Set \(\rho(dx)=f(x)\,dx\).
Its zero-extended square root belongs to \(W^{1,2}(\R)\), and
\(x(\sqrt f)'\in L^2(\R)\).
For \(0<c<1\), set
\[
 T=c\oplus I,\qquad
 f_c(y)=c^{-1}f(y/c),\qquad \mu=\rho^{\otimes\N}.
\]
Then \(T-I\) has rank one and
\[
 \nu_T=(f_c(y)\,dy)\otimes\rho^{\otimes\{2,3,\ldots\}}\ll\mu,
 \qquad \mu\not\ll\nu_T.
\]
Indeed, \(f_c\) is positive on \((-ca,ca)\) and zero on the rest
of \((-a,a)\).
Nevertheless,
\[
 \Aff(\mu,\nu_T)=\int\sqrt{f(y)f_c(y)}\,dy>0,
 \qquad \sup_n E_n(T)<\infty.
\]
The affinity sees the overlap on \((-ca,ca)\). The inverse row condition
sees the missing mass on the rest of \((-a,a)\). The smooth boundary makes
the Hellinger cost small, but it cannot restore that mass.

\paragraph{Gamma scale changes and the effect of centering.}

The half-line classification concerns linear maps of centered coordinates.
A dilation about the endpoint of a Gamma support is an affine map after
centering. Its compensating translation keeps the endpoint fixed, allowing
the familiar family of scale changes to reappear.
For the uncentered shape-\(r_\Gamma\), scale-one Gamma law, direct integration
gives
\[
  \int_0^\infty\sqrt{f(y)c^{-1}f(y/c)}\,dy
 =\left(\frac{2\sqrt c}{1+c}\right)^{r_\Gamma}
 =\cosh\!\left(\frac{\log c}{2}\right)^{-r_\Gamma}.
\]
Since \(\log\cosh(u/2)\sim u^2/8\), Kakutani gives
\begin{equation}
 \bigotimes_k\Law(c_kY)\sim\bigotimes_k\Law(Y)
 \quad\Longleftrightarrow\quad
 \sum_k(\log c_k)^2<\infty,
 \qquad c_k>0.
 \label{eq:gamma-scale-product-equivalence}
\end{equation}
Equivalently, \(\sum_k(c_k-1)^2<\infty\).  For
\(X=(Y-r_\Gamma)/\sqrt{r_\Gamma}\), the dilation \(Y\mapsto cY\) becomes
\[
 X\longmapsto cX+\sqrt{r_\Gamma}(c-1),
\]
rather than the homogeneous map \(X\mapsto cX\).

\begin{corollary}[Coordinatewise affine equivalence for Gamma products]
Let \(\mu\) be the product of the shape-\(r_\Gamma\) Gamma law centered with
variance one.
For \(T\in\GL(H)\) and \(b\in\ell^2\), the law of
\(\Phi_T(x)+b\) under \(\mu\) is equivalent to \(\mu\) if and only if
there are a coordinate
permutation \(\pi\) and numbers \(c_i>0\) such that
\[
 (Tx)_i=c_ix_{\pi(i)},\qquad
 b_i=\sqrt{r_\Gamma}(c_i-1),\qquad
 \sum_i(\log c_i)^2<\infty.
\]
\end{corollary}
\begin{proof}
Set \(S=T^{-1}\). Cauchy--Schwarz and
\eqref{eq:measurable-composition} give the inverse affine map
\(z\mapsto\Phi_S(z)-Sb\). The identities apply to the affine orbit
because \(\ell^2\) translations preserve the bound
\eqref{eq:covariance-domination}. Conditioning finite row sums as in
\cref{thm:bounded-and-half-line-classification} shows that both \(T\) and
\(T^{-1}\) have nonnegative entries: a negative coefficient would create an
unbounded lower tail that no fixed shift can remove.  The inverse identities
then make \(T\) monomial with positive weights \(c_i\).  Equality of the
scalar lower endpoints forces
\(-\sqrt{r_\Gamma}c_i+b_i=-\sqrt{r_\Gamma}\).
Equation \eqref{eq:gamma-scale-product-equivalence} completes both directions.
\end{proof}

\paragraph{An asymmetric example with no fourth moment.}

The finite-variance classification does not rely on moment expansions
beyond order two. To see the range this permits, consider the density
proportional to
\[
 (1+y^2)^{-(\tau+1)/2}e^{\vartheta\arctan y},\qquad
 2<\tau\le4,\quad\vartheta\in\R\setminus\{0\}.
\]
Its pre-normalization score is \(((\tau+1)y-\vartheta)/(1+y^2)\); it and its
product with \(y\) are bounded, and these bounds persist under normalization.
The fourth moment is infinite, and the two tails have
unequal constants. After centering and normalization, the main theorem
therefore classifies every bounded invertible mixing of this product.
The bounded Fourier tests still detect its coordinate structure, and the
tail exchanges still force exactly the Hilbert--Schmidt amount of deviation
from an ordinary permutation.
